\section{Prospective evaluation}
\subsection{Retrospective description versus pseudo-real-time analysis}
A full-sample historical factor provides a useful retrospective view
of financial disturbances, but its parameters and normalization may
use information from dates after the purported forecast origin.
Consequently, the primary forecasting experiment will instead use
historically available data and refit transformations, the index, the
trend model, and forecast calibration at each origin.
Retrospective and pseudo-real-time estimates will be labeled separately.

\subsection{Countries, events and benchmark forecasting methods}
Proposed countries comprise Mexico, the United States, the United
Kingdom, Japan and a euro-area aggregate. Predeclared case studies may
include the 2008--2009 financial crisis, European sovereign debt stress,
the COVID-19 shock, and later monetary and banking disturbances.
The 1994--1995 Mexican crisis is a separate possible extension because
the proposed contemporary indicators have limited older coverage.
All event and quiet-control windows will be fixed before confirmatory tests.

\subsection{Statistical evaluation}
For each index, assess matched-origin forecast RMSE, MAE, predictive
interval coverage and standardized exceedance rates against a random walk,
an autoregression, fixed trend smoothness, and (where meaningful) a
factor-transition forecast. Cross-index direct performance comparisons
require a predeclared \emph{common observable target}: indices are not
assumed to be interchangeable random variables. Dependence across
overlapping forecast horizons and shared country shocks must enter the
uncertainty analysis. The complete prospective protocol is EXP-0001.
